\documentclass[12pt,a4paper]{article}
\usepackage{luatexja}
\usepackage{luatexja-fontspec}
\usepackage{geometry}
\usepackage{listings}
\usepackage{xcolor}
\usepackage{booktabs}
\usepackage{tcolorbox}
\tcbuselibrary{skins}

\geometry{margin=25mm}

\definecolor{codebg}{HTML}{0A1520}
\definecolor{rusttext}{HTML}{D4DDE8}
\definecolor{codered}{HTML}{E05C2A}
\definecolor{codeblue}{HTML}{5BB0C0}
\definecolor{codegreen}{HTML}{A3D977}
\definecolor{codegray}{HTML}{4A6278}

\lstdefinelanguage{Rust}{
  keywords={fn, let, mut, if, else, loop, while, for, in, match, break, return, true, false, as},
  keywordstyle=\color{codered}\bfseries,
  ndkeywords={bool, i32, i64, u8, u32, f32, f64, usize, str},
  ndkeywordstyle=\color{codeblue},
  sensitive=true,
  morecomment=[l]{//},
  commentstyle=\color{codegray}\itshape,
  morestring=[b]",
  stringstyle=\color{codegreen},
}

\lstset{
  language=Rust,
  backgroundcolor=\color{codebg},
  basicstyle=\ttfamily\small\color{rusttext},
  frame=single,
  framerule=0.4pt,
  rulecolor=\color{codeblue!40},
  numbers=left,
  numberstyle=\tiny\color{codegray},
  numbersep=8pt,
  breaklines=true,
  showstringspaces=false,
  tabsize=4,
  xleftmargin=16pt,
}

\newtcolorbox{notebox}{
  colback=black!30,
  colframe=codeblue!60,
  leftrule=3pt, toprule=0.4pt, bottomrule=0.4pt, rightrule=0.4pt,
  fontupper=\small,
}

\newtcolorbox{pointbox}{
  colback=black!20,
  colframe=codered!60,
  leftrule=3pt, toprule=0.4pt, bottomrule=0.4pt, rightrule=0.4pt,
  fontupper=\small,
}

\begin{document}

\begin{center}
  {\small\ttfamily\color{codered} CHAPTER 02 · Tour of Rust}\\[6pt]
  {\LARGE\bfseries 第2章：基本制御フロー}\\[4pt]
  {\small\color{gray} 条件分岐・ループ・パターンマッチ}
\end{center}

\vspace{2mm}\hrule\vspace{6mm}

% ═══════════════════════════
\section{\texttt{if} / \texttt{else if} / \texttt{else}}

他の言語と違い、条件式に括弧 \texttt{()} が不要。比較演算子（\texttt{==}, \texttt{!=}, \texttt{<}, \texttt{>}, \texttt{<=}, \texttt{>=}）と論理演算子（\texttt{!}, \texttt{||}, \texttt{\&\&}）が使える。

比較演算子も演算の一種なので、\textbf{型が一致していないとコンパイルエラー}になる。例えば \texttt{f64} の変数と整数リテラル \texttt{42} を比較しようとするとエラーになるため、\texttt{42.0} と書いて型を合わせる必要がある。

\begin{lstlisting}
fn main() {
    let x = 41.9999_f64;
    if x < 42.0 {           // 42 ではなく 42.0 と書く
        println!("42より小さい");
    } else if x == 42.0 {
        println!("42に等しい");
    } else {
        println!("42より大きい");
    }
}
\end{lstlisting}

% ═══════════════════════════
\section{\texttt{loop}}

無限ループ。\texttt{break} で抜け出す。条件を書き忘れると整数がオーバーフローしてパニックするので注意。

\begin{lstlisting}
fn main() {
    let mut x = 0;
    loop {
        x += 1;
        if x == 42 {
            break;
        }
    }
    println!("{}", x);  // 42
}
\end{lstlisting}

% ═══════════════════════════
\section{\texttt{while}}

条件が \texttt{true} の間ループし続ける。条件が \texttt{false} になったら終了。出力がなくてもエラーではない（\texttt{println!} を書かなければ何も表示されないだけ）。

\begin{lstlisting}
fn main() {
    let mut x = 0;
    while x != 42 {
        x += 1;
    }
    println!("{}", x);  // 42
}
\end{lstlisting}

% ═══════════════════════════
\section{\texttt{for}}

イテレータを使って繰り返す。Rustでは \texttt{..} を使うだけで整数のシーケンス（イテレータ）を簡単に作れる。

\begin{lstlisting}
fn main() {
    // 0..5 → 0, 1, 2, 3, 4（終端を含まない）
    for x in 0..5 {
        println!("{}", x);
    }

    // 0..=5 → 0, 1, 2, 3, 4, 5（終端を含む）
    for x in 0..=5 {
        println!("{}", x);
    }

    // 配列を直接回す
    let nums = [10, 20, 30];
    for n in nums {
        println!("{}", n);
    }
}
\end{lstlisting}

\begin{notebox}
\texttt{..} と \texttt{..=} の違いは終端を含むかどうかだけ。Pythonの \texttt{range(0, 5)} が \texttt{0..5} に相当する。
\end{notebox}

% ═══════════════════════════
\section{\texttt{match}}

値が何かによって分岐する。他言語の \texttt{switch} に相当するが、\textbf{全パターンを網羅しないとコンパイルエラー}になる点が大きな違い。

\texttt{match} の構造は「パターン \texttt{=>} 処理」で、\texttt{=>} は矢印（不等号ではない）。

\begin{lstlisting}
fn main() {
    let x = 42;

    match x {
        0 => {
            println!("ゼロ");
        }
        1 | 2 => {              // | は「または」（1か2のどちらか）
            println!("1か2");
        }
        3..=9 => {              // 範囲指定（3以上9以下）
            println!("3から9");
        }
        matched_num @ 10..=100 => {  // @ でマッチした値に名前をつける
            println!("{}は10から100の間", matched_num);
        }
        _ => {                  // _ はワイルドカード（それ以外全部）
            println!("それ以外");
        }
    }
}
\end{lstlisting}

\begin{notebox}
\texttt{1 | 2} は「1または2」という列挙。\texttt{1..=2}（範囲）とは書き方が違うが、今回の場合は結果は同じ。\\
\texttt{matched\_num @ 10..=100} は「10以上100以下にマッチして、その値を \texttt{matched\_num} という名前で使う」という意味。
\end{notebox}

% ═══════════════════════════
\section{\texttt{loop} から値を返す}

\texttt{break} に値を付けると、\texttt{loop} 全体がその値を返す式になる。

\begin{lstlisting}
fn main() {
    let mut x = 0;
    let result = loop {
        x += 1;
        if x == 42 {
            break x * 2;    // break に値を付けて返す
        }
    };
    println!("{}", result); // 84
}
\end{lstlisting}

% ═══════════════════════════
\section{ブロック式から値を返す}

Rustでは \texttt{if}、\texttt{match}、\texttt{\{ \}} ブロックはすべて値を返す\textbf{式}。最後の行にセミコロンを付けなければその値が返る（第1章の関数と同じルール）。

これを使うと、Pythonの三項演算子（\texttt{-1 if x < 42 else 1}）と同じことが \texttt{if} 式で書ける。

\begin{lstlisting}
fn main() {
    let x = 42;

    // if を三項演算子のように使う（Pythonの「A if 条件 else B」と同じ発想）
    let label = if x < 42 { "小さい" } else { "大きい" };
    println!("{}", label);

    // ブロックも値を返せる
    let msg = {
        let y = x * 2;
        format!("x={}, 2倍={}", x, y)  // セミコロンなし → これが返り値
    };
    println!("{}", msg);
}
\end{lstlisting}

\begin{pointbox}
\textbf{文と式の区別}\\
セミコロンあり → 文（値を返さない）。セミコロンなし → 式（値を返す）。この区別がRust全体で一貫して使われる重要なルール。
\end{pointbox}

% ═══════════════════════════
\section*{まとめ}

\begin{center}
\begin{tabular}{ll}
\toprule
\textbf{構文} & \textbf{説明} \\
\midrule
\texttt{if x < 42 \{ \}} & 条件分岐（括弧不要）\\
\texttt{loop \{ break; \}} & 無限ループ + break \\
\texttt{while x != 0 \{ \}} & 条件付きループ \\
\texttt{for x in 0..5 \{ \}} & イテレータループ（終端含まず）\\
\texttt{for x in 0..=5 \{ \}} & イテレータループ（終端含む）\\
\texttt{match x \{ \_ => \}} & 網羅的パターンマッチ \\
\texttt{break x * 2} & loopから値を返す \\
\texttt{let v = if … \{ \} else \{ \}} & if式（三項演算子代わり）\\
\bottomrule
\end{tabular}
\end{center}

\vspace{6mm}
\hrule
\vspace{2mm}
{\small\color{gray}\textit{Tour of Rust — 第2章 / 基本制御フロー}}

\end{document}
